\begin{titlepage}
\begin{adjustwidth}{6mm}{6mm}

\makebox[30mm]{\vspace{20mm}}

\begin{center}
\zihao{-2}{\heiti 研究生学位论文的独创性声明}
\end{center}
\vspace{5mm}
\begin{spacing}{1.5}
\par {\zihao{4}\STZhongsong  本人声明，所呈交的论文是本人在导师指导下进行的研究工作及取得的研究成果。尽我所知，除了文中特别加以标注和致谢的地方外，论文中不包含其他人已经发表或撰写过的研究成果，也不包含为获得武汉理工大学或其他教育机构的学位或证书而使用过的材料。与我一同工作的同志对本研究所做的任何贡献均已在论文中作了明确的说明并表示了谢意。}
\end{spacing}
\vspace{3mm}
\hspace{40mm}{\zihao{4}\STZhongsong 签\hspace{3mm}名：} \CJKunderline{\makebox[30mm]{}} \hspace{8mm} {\zihao{4}\STZhongsong 日\hspace{3mm}期：} \CJKunderline{\makebox[30mm]{}}


\vspace{40mm}


\begin{center}
\zihao{-2}\heiti 学位论文使用授权书
\end{center}
\vspace{3mm}
\begin{spacing}{1.5}
\par {\zihao{4}\STZhongsong 本人完全了解武汉理工大学有关保留、使用学位论文的规定，即学校有权保留并向国家有关部门或机构送交论文的复印件和电子版，允许论文被查阅和借阅。本人授权武汉理工大学可以将本学位论文的全部内容编入有关数据库进行检索，可以采用影印、缩印或其他复制手段保存或汇编本学位论文。同时授权经武汉理工大学认可的国家有关机构或论文数据库使用或收录本学位论文,并向社会公众提供信息服务。}
\end{spacing}
\vspace{3mm}
\par {\zihao{4}\STZhongsong（保密的论文在解密后应遵守此规定）}

\vspace{20mm}
\par {\zihao{4}\STZhongsong 研究生（签名）：} \hspace{20mm} {\zihao{4}\STZhongsong 导师（签名）：} \hspace{20mm} {\zihao{4}\STZhongsong 日期：}

\end{adjustwidth}
\end{titlepage}


\begin{titlepage}
\begin{adjustwidth}{6mm}{6mm}

\begin{center}
\zihao{-2}\heiti 学位申请人AI工具使用声明
\end{center}
\vspace{3mm}
\begin{spacing}{1.5}
\par {\zihao{4}\STZhongsong 本人声明：在本篇学位论文或实践成果总结报告中，作者使用了［填写AI工具名称］进行［填入具体用途，举例：辅助文献检索与整理/图表和音视频辅助制作/代码调试与分析辅助/语法校对与格式检查等］，所有AI辅助生成合成内容均已标注。学位论文或实践成果创新观点、研究数据、实验实践设计、核心结论均由本人独立完成。本人对使用该工具辅助完成的内容进行了合理编辑并进行了真实性、合法性审查，对本篇学位论文或实践成果总结报告的内容承担全部法律和学术伦理责任。}
\end{spacing}
\vspace{3mm}
\par {\zihao{4}\STZhongsong（本人声明：在本篇学位论文或实践成果总结报告中，未使用AI工具。）}

\vspace{15mm}
\par {\zihao{4}\STZhongsong 本人签名：} \hspace{30mm} {\zihao{4}\STZhongsong 日期：}

\vspace{20mm}
\par {\zihao{4}\STZhongsong 导师意见：} 
\par {\zihao{4}\STZhongsong ［填写对学位申请人AI工具使用声明的审核意见］}

\end{adjustwidth}
\end{titlepage}


